\documentclass{article}
\usepackage[utf8]{inputenc}
\usepackage[T1]{fontenc}
\usepackage{geometry}
\geometry{
    a4paper,
    left=2.5cm,
    right=2.5cm,
    top=2.5cm,
    bottom=2.5cm
}
\usepackage{graphicx}
\usepackage{float}
\usepackage{hyperref}
\usepackage[english]{babel}
\usepackage{amsmath}
\usepackage{amssymb}
\usepackage{array}
\usepackage{booktabs}
\usepackage{caption}
\usepackage{subcaption}
\usepackage{enumitem}
\usepackage{fancyhdr}
\usepackage{titlesec}
\usepackage{url}
\usepackage{textcomp}
\usepackage{xcolor}
\usepackage{listings}
\usepackage{multirow}
\usepackage{longtable}
\usepackage{tabularx}
\usepackage{times}

% Set up page style
\pagestyle{fancy}
\fancyhf{}
\rfoot{\thepage}
\renewcommand{\headrulewidth}{0pt}

% Title formatting
\titleformat{\chapter}[display]{\normalfont\bfseries}{\chaptertitle\thechapter}{0pt}{\Huge}
\titlespacing*{\chapter}{0pt}{50pt}{40pt}

% Document info
\title{Template Document}
\author{Academic Support Agent}
\date{\today}

\begin{document}
\maketitle

\begin{frame}[plain]
\maketitle
\end{frame}

\section[N zov sekcie 1]{Dlh  n zov sekcie 1}
\subsection[N zov podsekcie 1]{Dlh  n zov podsekcie 1}

\begin{frame}{Nadpis}{Podnadpis}
oby ajn  text, \structure{ trukt ra str nky}, \alert{zv raznen  text}
\begin{itemize}
  \item polo ka odr kov ho zoznamu na jeden riadok
  \item polo ka odr kov ho zoznamu dlh , pr li  dlh  (aby sa zalomila), ktor 
    obsahuje aj \alert{zv raznen  text}
  \begin{itemize}
    \item odr ka druhej  rovne
    \begin{itemize}
      \item odr ka tretej  rovne
    \end{itemize}
    \item \alert{zv raznen  odr ka druhej  rovne}
  \end{itemize}
\end{itemize}
\begin{enumerate}
  \item a  slovan  odr ka
  \begin{enumerate}
    \item odr ka druhej  rovne se vzor kom
      \[ E = mc^2 \]
      a s odkazom na bibliografick  cit ciu \cite{einstein1905tragheit}
  \end{enumerate}
\end{enumerate}
\end{frame}

\subsection[N zov podsekcie 2]{Dlh  n zov podsekcie 2}

\begin{frame}{Textov  bloky}
text nad blokom\footnote{text pozn mky s \url{https://adresou.sk}}
\begin{block}{Blok}
  text v bloku
\end{block}
\begin{exampleblock}{Blok s pr kladom}
  text v bloku
\end{exampleblock}
\begin{alertblock}{Zv raznen  blok}
  text v bloku
\end{alertblock}
\end{frame}

\begin{frame}{Obr zky}
\begin{figure}
  \includegraphics[width=.5\textwidth,height=.5\textheight,keepaspectratio]{cow-black.mps}
  \caption{Krava hol t nska v stoji%
    \footnote{%
      Kedy  e si t  krava znova  ahne?
      Odpove  v s prekvap .
      \cite{tolkamp10cows}
    }%
  }
\end{figure}
\end{frame}

\subsection[N zov podsekcie 3]{Dlh  n zov podsekcie 3}

\def\age(#1-#2-#3){%
  \the\numexpr(
    \year - #1
    \ifnum\month<#2
      - 1
    \else
      \ifnum\month=#2
        \ifnum\day<#3
          - 1
        \fi
      \fi
    \fi
  )%
}

\begin{frame}{Tabu ky}
\begin{table}
  \begin{tabular}{llc}
    Meno & Priezvisko & Vek \\ \midrule
    Albert & Einstein & \age(1879-03-14) \\
    Marie & Curie & \age(1867-11-07) \\
    Thomas & Edison & \age(1847-02-11) \\
  \end{tabular}
  \caption{Ve k  vedci z 19. storo ia}
\end{table}
\end{frame}

\section{\bibname}
\begin{frame}[t, allowframebreaks]{\bibname}
\printbibliography[heading=none]
\end{frame}

\begin{frame}[plain]
\vfill
\centerline{ akujem v m za pozornos !}
\vfill\vfill
\end{frame}

\makeoutro
\end{document}